%==============================================================================
% jarvis-dirtree.tex
% Pretty directory trees for the manual, built on dirtree + fontawesome5.
%
% pdfLaTeX cannot render colour emoji (those need LuaLaTeX/XeLaTeX), so we use
% FontAwesome 5 glyphs coloured in the Jarvis palette as folder/file icons.
%
% ICON STYLE: outline only ("regular"), never filled ("solid").  Every icon below
% uses the [regular] variant.  Note that FontAwesome 5 *Free* ships only a limited
% regular set: file-csv and file-contract are solid-only, so CSV uses the outline
% spreadsheet glyph and there is no contract icon here.
%
% Usage:
%   \begin{jfiletree}
%   \dirtree{%
%     .1 \dtfold{MyScan/}.
%     .2 \dtfold{bin/}.
%     .3 \dtyaml{task.yaml}.
%     .2 \dtmark{.jarvis-project.json}\DTcomment{project marker}.
%   }
%   \end{jfiletree}
%
% \DTcomment{...} adds a right-aligned, dotted-leader annotation.
%==============================================================================

\usepackage{dirtree}
\usepackage{fontawesome5}

% Icon colours (Jarvis palette).
\definecolor{jfolder}{HTML}{75A3FA}   % folders: Jarvis blue
\definecolor{jfileicon}{HTML}{7B8794} % files: neutral slate
\definecolor{jmarkicon}{HTML}{B3261E} % project markers: Jarvis red

% Generic builder: an outline (regular) icon in a colour, then the monospace name.
\newcommand{\dticon}[2]{\textcolor{#1}{#2}\hspace{0.45em}}

% --- Folders (Jarvis blue) ----------------------------------------------------
\newcommand{\dtfold}[1]{\dticon{jfolder}{\faFolder[regular]}\texttt{#1}}
\newcommand{\dtfoldopen}[1]{\dticon{jfolder}{\faFolderOpen[regular]}\texttt{#1}}

% --- Typed files (neutral slate) ----------------------------------------------
\newcommand{\dtfile}[1]{\dticon{jfileicon}{\faFile[regular]}\texttt{#1}}        % generic file
\newcommand{\dtcode}[1]{\dticon{jfileicon}{\faFileCode[regular]}\texttt{#1}}     % yaml/json/config/script
\newcommand{\dtjson}[1]{\dticon{jfileicon}{\faFileCode[regular]}\texttt{#1}}     % json (alias of \dtcode)
\newcommand{\dtcsv}[1]{\dticon{jfileicon}{\faFileExcel[regular]}\texttt{#1}}     % csv / tabular data
\newcommand{\dtdata}[1]{\dticon{jfileicon}{\faChartBar[regular]}\texttt{#1}}     % hdf5 / results data
\newcommand{\dtdoc}[1]{\dticon{jfileicon}{\faFile[regular]}\texttt{#1}}          % readme / markdown / text
\newcommand{\dtimage}[1]{\dticon{jfileicon}{\faFileImage[regular]}\texttt{#1}}   % image / plot
\newcommand{\dtpdf}[1]{\dticon{jfileicon}{\faFilePdf[regular]}\texttt{#1}}       % pdf
\newcommand{\dtarchive}[1]{\dticon{jfileicon}{\faFileArchive[regular]}\texttt{#1}} % tar.gz / zip
\newcommand{\dtnote}[1]{\dticon{jfileicon}{\faStickyNote[regular]}\texttt{#1}}   % note

% --- Project markers (Jarvis red) ---------------------------------------------
\newcommand{\dtmark}[1]{\dticon{jmarkicon}{\faFile[regular]}\texttt{#1}}         % .jarvis-project.json
\newcommand{\dtyaml}[1]{\dticon{jmarkicon}{\faFileCode[regular]}\texttt{#1}}     % jarvis.project.yaml marker

% Width-constrained wrapper so dirtree lays out inside the text column (its
% dotted \DTcomment leaders run to the column's right edge).
\newenvironment{jfiletree}
  {\par\addvspace{0.6\baselineskip}\noindent\begin{minipage}{\linewidth}\footnotesize}
  {\end{minipage}\par\addvspace{0.6\baselineskip}}

%==============================================================================
% Inline file/folder mentions.
% \file{...} (defined plain in jarvis-preamble.tex) is upgraded here to prepend
% the matching outline icon, so EVERY inline file/folder reference in the prose
% uses the same icon set as the directory trees.  The icon is chosen from the
% name: a trailing "/" or ">" (e.g. a "<uuid>" placeholder dir) => folder;
% otherwise by extension; anything else => a generic file.
%==============================================================================
\usepackage{xstring}

% half-space before the icon + icon (coloured) + half-space + path name (dark green)
% (\, is a thin space, ~half a normal interword space)
\newcommand{\jfilenode}[3]{\,\textcolor{#1}{#2}\,\textcolor{jPath}{\jtt{#3}}}

\renewcommand{\file}[1]{%
  \IfEndWith{#1}{/}{\jfilenode{jfolder}{\faFolder[regular]}{#1}}{%
  \IfEndWith{#1}{>}{\jfilenode{jfolder}{\faFolder[regular]}{#1}}{%
  \IfEndWith{#1}{.yaml}{\jfilenode{jfileicon}{\faFileCode[regular]}{#1}}{%
  \IfEndWith{#1}{.json}{\jfilenode{jfileicon}{\faFileCode[regular]}{#1}}{%
  \IfEndWith{#1}{.csv}{\jfilenode{jfileicon}{\faFileExcel[regular]}{#1}}{%
  \IfEndWith{#1}{.hdf5}{\jfilenode{jfileicon}{\faChartBar[regular]}{#1}}{%
  \IfEndWith{#1}{.md}{\jfilenode{jfileicon}{\faFile[regular]}{#1}}{%
  \IfEndWith{#1}{.txt}{\jfilenode{jfileicon}{\faFile[regular]}{#1}}{%
  \IfEndWith{#1}{.png}{\jfilenode{jfileicon}{\faFileImage[regular]}{#1}}{%
  \IfEndWith{#1}{.pdf}{\jfilenode{jfileicon}{\faFilePdf[regular]}{#1}}{%
  \jfilenode{jfileicon}{\faFile[regular]}{#1}}}}}}}}}}}%
}
